%! Probability Models — Unit 3, Lesson 3.2 — Student Packet Cover Sheet
% THREE packet rows — Warm-Up, Guided Notes & Practice, Homework. The homework is the individual
% practice, assigned at Close & Assign and done outside class, and it is scored 1 to 5. The cover
% is a contents list, with no Score column.
\documentclass[12pt]{article}
\usepackage{saar-article}
\usepackage{saar-boxes}
\usepackage{ltablex}
\keepXColumns
% Measured banner: the band grows to fit the title block (copied verbatim from the Unit 2 covers).
\newsavebox{\bannerbox}
\newlength{\bannerht}
\newlength{\coverbannerpad}
\setlength{\coverbannerpad}{0.16in}       % breathing room above and below the text
\newcommand{\coverbanner}[2]{%
  \sbox{\bannerbox}{%
    \begin{minipage}{\dimexpr\paperwidth-1.4in\relax}%
      \centering
      {\color{white}\textbf{\LARGE \CourseName}}\\[4pt]
      {\color{frostmid}\large\textbf{#1}}\\[6pt]
      {\color{white}\textbf{\large #2}}%
    \end{minipage}}%
  \setlength{\bannerht}{\dimexpr\ht\bannerbox+\dp\bannerbox+2\coverbannerpad\relax}%
  \noindent\begin{tikzpicture}[remember picture, overlay]
    \fill[cerulean] (current page.north west)
      rectangle ([yshift=-\bannerht]current page.north east);
    \node[anchor=north, inner sep=0pt]
      at ([yshift=-\coverbannerpad]current page.north) {\usebox{\bannerbox}};
  \end{tikzpicture}\par
  % The text body starts 0.6in below the page edge (geometry top=0.6in), so reserve only what
  % the band overruns that, plus a gap under the band.
  \vspace*{\dimexpr\bannerht-0.6in+0.16in\relax}%
}

\begin{document}

% Full-bleed cerulean banner, sized to the title block.
\coverbanner{Unit 3}{Lesson 3.2 \quad Probability Models}

\vspace{0.06in}
% NAMESTRIP: this is the ONLY component that carries the name row.
\namedateperiod
\vspace{0.18in}

% The standard codes stay in the teacher-facing plan.
\begin{learningtargetbox}
\small
\begin{enumerate}[leftmargin=1.5em, itemsep=5pt]
  \item \ldots{} list a \textbf{sample space} and an \textbf{event}, find a theoretical
        probability, and use it to predict an expected count.
  \item \ldots{} use the \textbf{complement} rule, and tell \textbf{complementary} events from
        \textbf{mutually exclusive} events, in context.
  \item \ldots{} organize one situation as a tree diagram and as a two-way table, and read
        probabilities from each.
  \item \ldots{} read a Venn diagram for the intersection, the union, ``neither,'' and ``only''
        --- and count ``or'' without counting the overlap twice.
\end{enumerate}
\end{learningtargetbox}

\vspace{0.14in}

\begin{tocbox}
\small
\renewcommand{\arraystretch}{1.5}
\begin{tabularx}{\linewidth}{@{} c l X @{}}
\rowcolor{frost}
\textbf{\#} & \textbf{Component} & \textbf{Description} \\
\midrule
1 & Warm-Up & Simplifying fractions and moving among fractions, decimals, and percents, with the
             numbers the notes use \\
2 & Guided Notes \& Practice & A food-truck festival: a prize wheel, complements, combo orders on
             a tree and in a table, and a survey on a Venn diagram --- then a sports roster we
             work together \\
3 & Homework & A bookstore loyalty program: a rewards wheel and a members survey.
             \textbf{Scored 1 to 5, assigned today, due the first class after two study
             halls.} \\
\end{tabularx}
\end{tocbox}

\vspace{0.14in}

% Keep in Mind carries the lesson's CONTENT — the rule, the test, the trap — never its process.
\begin{remindbox}
\small
With equally likely outcomes, $P(\text{event})$ is the outcomes in the event over the outcomes in
the \textbf{sample space}. The \textbf{complement} rule is $P(\text{not } A) = 1 - P(A)$.
\textbf{Mutually exclusive} events cannot both happen; \textbf{complementary} events cannot both
happen \emph{and} together cover everything. \textbf{Watch out:} a whole circle is not the
``only'' region --- take out the overlap. And adding two circle totals counts the overlap twice:
to count ``or,'' add the regions inside the circles, each one once.
\end{remindbox}

\end{document}
